\documentclass[12pt]{article}
\usepackage[a4paper,margin=1in]{geometry}
\usepackage{iftex}
\ifXeTeX
\usepackage{fontspec}
\defaultfontfeatures{Ligatures=TeX}
\setmainfont{Times New Roman}
\else
\usepackage[T1]{fontenc}
\usepackage[utf8]{inputenc}
\usepackage{newtxtext}
\usepackage{newtxmath}
\fi
\usepackage{enumitem}
\usepackage{hyperref}
\usepackage{parskip}
\usepackage{microtype}
\usepackage[normalem]{ulem}
\usepackage{xcolor}

\newcommand{\practiceid}[1]{\texttt{#1}}
\newlength{\readinggutter}
\setlength{\readinggutter}{2.8em}
\newlength{\readingfirstindent}
\setlength{\readingfirstindent}{1.5em}
\newlength{\questionblockgap}
\setlength{\questionblockgap}{0.45\baselineskip}
\newlength{\exampleblockgap}
\setlength{\exampleblockgap}{0.65\baselineskip}
\newlength{\passageparagraphgap}
\setlength{\passageparagraphgap}{0.45\baselineskip}
\newenvironment{exampleblock}{\par\vspace{0.35\baselineskip}\begingroup\raggedright\setlength{\parskip}{0pt}\setlength{\parindent}{0pt}}{\par\endgroup\vspace{\exampleblockgap}}
\newcommand{\readinglineindent}[2]{%
\par\noindent%
\hangindent=\dimexpr\readinggutter+0.9em\relax%
\hangafter=1%
\makebox[\readinggutter][r]{\scriptsize\color{gray}#1}\hspace{0.9em}\hspace*{\readingfirstindent}#2%
\par}
\newcommand{\readingline}[2]{%
\par\noindent%
\hangindent=\dimexpr\readinggutter+0.9em\relax%
\hangafter=1%
\makebox[\readinggutter][r]{\scriptsize\color{gray}#1}\hspace{0.9em}#2%
\par}

\setlength{\parindent}{0pt}
\setlength{\parskip}{0.5em}
\setlist[enumerate]{topsep=0.25em,itemsep=0.18em,parsep=0pt,partopsep=0pt}
\setlength{\emergencystretch}{2em}

\begin{document}
\section*{TOEFL Practice 04 - Tryout}
Source of truth: DOCX only.
\section*{Listening Comprehension}
\subsection*{Part A}
DIRECTIONS: In Part A you will hear short conversations between two speakers. At the end of each conversation, a third speaker will ask a question about what the first two speakers said. Each conversation and each question will be spoken only one time. Therefore, you must listen carefully to understand what each speaker says. After you hear a conversation and the question, read the four choices and select the one that is the best answer to the question the speaker asked. Then, on your answer. sheet, find the number of the question and blacken the space that corresponds to the letter for the answer you have chosen. Blacken the space completely so that the letter inside the space does not show.\par
\begin{exampleblock}
Listen to the following example.\par
On the recording, you hear:\par
(Man) Does the car need to be filled?\par
(Woman) Mary stopped at the gas station on her way home.\par
(Narrator) What does the woman mean?\par
In your test book, you will read:\par
(A) Mary bought some food.\par
(B) Mary had car trouble.\par
(C) Mary went shopping.\par
(D) Mary bought some gas.\par
From the conversation you learn that Mary stopped at the gas station on her way home. The best answer to the question "Does the car need to be filled?" is (D), "Mary bought some gas." Therefore, the correct answer is (D).\par
Now let us begin Part A with question number 1.\par
\end{exampleblock}
\begin{exampleblock}
\textbf{Transkrip rekonstruksi AI (draft; bukan resmi)}\par
Man: I heard Professor White's advanced seminar is one of the most demanding courses in the department.\par
Woman: It is, but his students always speak of him with real admiration.\par
Narrator: What does the woman imply?\par
\end{exampleblock}
\noindent\textbf{1.}
\begin{enumerate}[label=(\Alph*),leftmargin=*,itemsep=0.2em]
\item Professor White holds advanced classes.
\item Professor White's students have graduated.
\item Students have great respect for Professor White.
\item Students are taking Professor White's seminar.
\end{enumerate}
\par
\vspace{0.25\baselineskip}
\begin{exampleblock}
\textbf{Transkrip rekonstruksi AI (draft; bukan resmi)}\par
Man: That was an excellent dinner.\par
Woman: I'm glad you enjoyed it. Would you care for some dessert?\par
Narrator: What does the woman mean?\par
\end{exampleblock}
\noindent\textbf{2.}
\begin{enumerate}[label=(\Alph*),leftmargin=*,itemsep=0.2em]
\item You have to be careful in the desert.
\item You shouldn't eat so many sweets.
\item May I offer you some dessert?
\item Can you help me with the sweets?
\end{enumerate}
\par
\vspace{0.25\baselineskip}
\begin{exampleblock}
\textbf{Transkrip rekonstruksi AI (draft; bukan resmi)}\par
Woman: Could you lend me a little money for the vending machine?\par
Man: Sorry, I left my wallet at home.\par
Narrator: What does the man mean?\par
\end{exampleblock}
\noindent\textbf{3.}
\begin{enumerate}[label=(\Alph*),leftmargin=*,itemsep=0.2em]
\item He never has any money.
\item He has no money on him.
\item He is out getting cash.
\item He has been out of work.
\end{enumerate}
\par
\vspace{0.25\baselineskip}
\begin{exampleblock}
\textbf{Transkrip rekonstruksi AI (draft; bukan resmi)}\par
Man: Should I leave the copying machine on after I finish these copies?\par
Woman: No, please turn it off when you're done.\par
Narrator: What does the woman mean?\par
\end{exampleblock}
\noindent\textbf{4.}
\begin{enumerate}[label=(\Alph*),leftmargin=*,itemsep=0.2em]
\item The machine shouldn't be left running.
\item Turn on the copying machine when you leave.
\item The machine has to be turned around.
\item You made enough copies already.
\end{enumerate}
\par
\vspace{0.25\baselineskip}
\begin{exampleblock}
\textbf{Transkrip rekonstruksi AI (draft; bukan resmi)}\par
Man: Did you hear that Tom passed his Latin course?\par
Woman: Really? I thought everyone expected him to fail it.\par
Narrator: What does the woman imply?\par
\end{exampleblock}
\noindent\textbf{5.}
\begin{enumerate}[label=(\Alph*),leftmargin=*,itemsep=0.2em]
\item Tom passed up a Latin course.
\item Tom was expected to fail Latin.
\item He's surprised Tom is late.
\item Tom's Latin class is last.
\end{enumerate}
\par
\vspace{0.25\baselineskip}
\begin{exampleblock}
\textbf{Transkrip rekonstruksi AI (draft; bukan resmi)}\par
Man: Did she spend one hundred fifteen dollars on that dress?\par
Woman: The dress was one hundred fifteen, and the matching jacket was ninety-five.\par
Narrator: How much did she spend?\par
\end{exampleblock}
\noindent\textbf{6.}
\begin{enumerate}[label=(\Alph*),leftmargin=*,itemsep=0.2em]
\item She buys exotic clothes.
\item She bought a \$115 dress.
\item She spent \$210.
\item She has a new house.
\end{enumerate}
\par
\vspace{0.25\baselineskip}
\begin{exampleblock}
\textbf{Transkrip rekonstruksi AI (draft; bukan resmi)}\par
Man: What do you think of the new engineering building?\par
Woman: To tell the truth, it isn't at all what I expected.\par
Narrator: What does the woman mean?\par
\end{exampleblock}
\noindent\textbf{7.}
\begin{enumerate}[label=(\Alph*),leftmargin=*,itemsep=0.2em]
\item The new engineering building is far from here.
\item The building is not what I thought it would be.
\item I don't know what I expected the engineer to do.
\item They are building a new house for the engineer.
\end{enumerate}
\par
\vspace{0.25\baselineskip}
\begin{exampleblock}
\textbf{Transkrip rekonstruksi AI (draft; bukan resmi)}\par
Woman: Have you finally made that appointment with the dentist?\par
Man: Not yet. I keep putting it off.\par
Narrator: What does the man mean?\par
\end{exampleblock}
\noindent\textbf{8.}
\begin{enumerate}[label=(\Alph*),leftmargin=*,itemsep=0.2em]
\item The man has put off going to a dentist.
\item The man likes getting his teeth cleaned.
\item She is not happy that the man went to a dentist.
\item Everyone should have their teeth cleaned regularly.
\end{enumerate}
\par
\vspace{0.25\baselineskip}
\begin{exampleblock}
\textbf{Transkrip rekonstruksi AI (draft; bukan resmi)}\par
Man: Do you think this hall will work for the wedding reception?\par
Woman: Not for two hundred guests. It won't hold nearly that many people.\par
Narrator: What does the woman imply?\par
\end{exampleblock}
\noindent\textbf{9.}
\begin{enumerate}[label=(\Alph*),leftmargin=*,itemsep=0.2em]
\item The rent for the hall is high.
\item The wedding hall is for rent.
\item The hall is too small.
\item There are not enough guests.
\end{enumerate}
\par
\vspace{0.25\baselineskip}
\begin{exampleblock}
\textbf{Transkrip rekonstruksi AI (draft; bukan resmi)}\par
Man: Could you cash this traveler's check for one thousand dollars?\par
Woman: I'm sorry. For checks that large, only the manager is authorized to cash them.\par
Narrator: What does the woman mean?\par
\end{exampleblock}
\noindent\textbf{10.}
\begin{enumerate}[label=(\Alph*),leftmargin=*,itemsep=0.2em]
\item She can cash the check.
\item Only the manager can cash the check.
\item Traveller's checks cannot be cashed there.
\item A thousand dollars is a large sum of money.
\end{enumerate}
\par
\vspace{0.25\baselineskip}
\begin{exampleblock}
\textbf{Transkrip rekonstruksi AI (draft; bukan resmi)}\par
Man: How did the housecleaning go yesterday?\par
Woman: Great. The crew washed all the windows and cleaned the rugs.\par
Narrator: What does the woman mean?\par
\end{exampleblock}
\noindent\textbf{11.}
\begin{enumerate}[label=(\Alph*),leftmargin=*,itemsep=0.2em]
\item They worked together as a team.
\item The windows and carpets are beautiful.
\item They now have clean windows and rugs.
\item They are always cleaning their house.
\end{enumerate}
\par
\vspace{0.25\baselineskip}
\begin{exampleblock}
\textbf{Transkrip rekonstruksi AI (draft; bukan resmi)}\par
Man: What was Linda's reaction when she won the award?\par
Woman: She was completely at a loss for words.\par
Narrator: What does the woman mean?\par
\end{exampleblock}
\noindent\textbf{12.}
\begin{enumerate}[label=(\Alph*),leftmargin=*,itemsep=0.2em]
\item She got lost one more time.
\item She didn't know what to say.
\item She is learning many new words.
\item She had less than the price of the book.
\end{enumerate}
\par
\vspace{0.25\baselineskip}
\begin{exampleblock}
\textbf{Transkrip rekonstruksi AI (draft; bukan resmi)}\par
Woman: Mark looks exhausted this morning.\par
Man: He really should stop keeping such late hours.\par
Narrator: What does the man mean about Mark?\par
\end{exampleblock}
\noindent\textbf{13.}
\begin{enumerate}[label=(\Alph*),leftmargin=*,itemsep=0.2em]
\item He is staying in this hotel.
\item He keeps going to bed late.
\item He lives by himself in this house.
\item He started early and finished late.
\end{enumerate}
\par
\vspace{0.25\baselineskip}
\begin{exampleblock}
\textbf{Transkrip rekonstruksi AI (draft; bukan resmi)}\par
Man: Is Ronald planning to travel this summer?\par
Woman: Yes, he says he wants to go abroad, but he hasn't decided which country yet.\par
Narrator: What does the woman mean?\par
\end{exampleblock}
\noindent\textbf{14.}
\begin{enumerate}[label=(\Alph*),leftmargin=*,itemsep=0.2em]
\item Ronald won't travel this summer.
\item She doesn't know to what country Ronald is going.
\item She wants to travel with Ronald.
\item Ronald wants to travel out of the country.
\end{enumerate}
\par
\vspace{0.25\baselineskip}
\begin{exampleblock}
\textbf{Transkrip rekonstruksi AI (draft; bukan resmi)}\par
Woman: The store said they couldn't deliver my dress today.\par
Man: Did you give them your address? They can't deliver it if they don't know where you live.\par
Narrator: What does the man imply?\par
\end{exampleblock}
\noindent\textbf{15.}
\begin{enumerate}[label=(\Alph*),leftmargin=*,itemsep=0.2em]
\item They can't deliver her dress.
\item They don't know where she lives.
\item They didn't touch her dress.
\item They aren't familiar with the city.
\end{enumerate}
\par
\vspace{0.25\baselineskip}
\begin{exampleblock}
\textbf{Transkrip rekonstruksi AI (draft; bukan resmi)}\par
Man: We still have forty miles to go, and the meeting starts in half an hour.\par
Woman: Then there is no way we'll get there on time.\par
Narrator: What does the woman mean?\par
\end{exampleblock}
\noindent\textbf{16.}
\begin{enumerate}[label=(\Alph*),leftmargin=*,itemsep=0.2em]
\item We'll be there on time.
\item We'll be late.
\item We'll take the next turn.
\item The turn is forty miles away.
\end{enumerate}
\par
\vspace{0.25\baselineskip}
\begin{exampleblock}
\textbf{Transkrip rekonstruksi AI (draft; bukan resmi)}\par
Man: The business school expected nine hundred new students this year.\par
Woman: But only about three hundred actually enrolled.\par
Narrator: What does the woman imply?\par
\end{exampleblock}
\noindent\textbf{17.}
\begin{enumerate}[label=(\Alph*),leftmargin=*,itemsep=0.2em]
\item Only a third of the students are enrolled in the business school.
\item The third enrollment prediction for the business school was accurate.
\item The enrollment is three times higher than what was predicted.
\item The enrollment figures are about 66 percent lower than expected.
\end{enumerate}
\par
\vspace{0.25\baselineskip}
\begin{exampleblock}
\textbf{Transkrip rekonstruksi AI (draft; bukan resmi)}\par
Man: Did you hear that Mr. Benson has started taking singing lessons?\par
Woman: At his age? That's certainly an unusual time to begin.\par
Narrator: What does the woman mean?\par
\end{exampleblock}
\noindent\textbf{18.}
\begin{enumerate}[label=(\Alph*),leftmargin=*,itemsep=0.2em]
\item He likes to sing old songs.
\item The rumors about him are unbelievable.
\item To start singing at his age is unusual.
\item You shouldn't talk about him this way.
\end{enumerate}
\par
\vspace{0.25\baselineskip}
\begin{exampleblock}
\textbf{Transkrip rekonstruksi AI (draft; bukan resmi)}\par
Man: I explained the assignment three times, but the students still asked the same questions afterward.\par
Woman: It sounds as if they aren't paying attention to what you say.\par
Narrator: What does the woman mean?\par
\end{exampleblock}
\noindent\textbf{19.}
\begin{enumerate}[label=(\Alph*),leftmargin=*,itemsep=0.2em]
\item He doesn't feel good, and he won't go to class.
\item He couldn't hear the lecture because students were talking.
\item Students don't pay attention to what he says.
\item In this class, the same thing happens every day.
\end{enumerate}
\par
\vspace{0.25\baselineskip}
\begin{exampleblock}
\textbf{Transkrip rekonstruksi AI (draft; bukan resmi)}\par
Man: What did you and your sister do after dinner?\par
Woman: We took a walk along the river and watched the sunset.\par
Narrator: What did the woman and her sister do?\par
\end{exampleblock}
\noindent\textbf{20.}
\begin{enumerate}[label=(\Alph*),leftmargin=*,itemsep=0.2em]
\item They passed the river.
\item They went for a walk.
\item They visited her sister.
\item They were looking at the sunset.
\end{enumerate}
\par
\vspace{0.25\baselineskip}
\begin{exampleblock}
\textbf{Transkrip rekonstruksi AI (draft; bukan resmi)}\par
Man: Should we go down to the office for tickets to the game?\par
Woman: Let's call first and see whether any are still available.\par
Narrator: What will they probably do first?\par
\end{exampleblock}
\noindent\textbf{21.}
\begin{enumerate}[label=(\Alph*),leftmargin=*,itemsep=0.2em]
\item Stand in line
\item Try to order tickets
\item Go to the game
\item Call their office
\end{enumerate}
\par
\vspace{0.25\baselineskip}
\begin{exampleblock}
\textbf{Transkrip rekonstruksi AI (draft; bukan resmi)}\par
Man: Since the phone is always busy, maybe we should get another line.\par
Woman: I don't think we can afford another monthly bill right now.\par
Narrator: What does the woman mean?\par
\end{exampleblock}
\noindent\textbf{22.}
\begin{enumerate}[label=(\Alph*),leftmargin=*,itemsep=0.2em]
\item She doesn't need another phone.
\item She is not as busy as it seems.
\item They cant afford another line.
\item They have to use the phone less.
\end{enumerate}
\par
\vspace{0.25\baselineskip}
\begin{exampleblock}
\textbf{Transkrip rekonstruksi AI (draft; bukan resmi)}\par
Woman: Since I already took math at another school, can I skip the placement test?\par
Man: No. All new students have to take the math placement test, even with transfer credit.\par
Narrator: What does the man mean?\par
\end{exampleblock}
\noindent\textbf{23.}
\begin{enumerate}[label=(\Alph*),leftmargin=*,itemsep=0.2em]
\item The man administers math placement tests.
\item The woman should take a math course.
\item The woman has to take the placement test.
\item The woman can transfer her math credit.
\end{enumerate}
\par
\vspace{0.25\baselineskip}
\begin{exampleblock}
\textbf{Transkrip rekonstruksi AI (draft; bukan resmi)}\par
Man: The roses need pruning, and the lawn has to be reseeded.\par
Woman: Then we should hire someone who really knows how to care for plants.\par
Narrator: Whom do they need?\par
\end{exampleblock}
\noindent\textbf{24.}
\begin{enumerate}[label=(\Alph*),leftmargin=*,itemsep=0.2em]
\item A florist
\item A gardener
\item A barber
\item A custodian
\end{enumerate}
\par
\vspace{0.25\baselineskip}
\begin{exampleblock}
\textbf{Transkrip rekonstruksi AI (draft; bukan resmi)}\par
Man: What are you doing in the garage?\par
Woman: I'm just listening to the final round of the contest on the radio.\par
Narrator: What is the woman doing?\par
\end{exampleblock}
\noindent\textbf{25.}
\begin{enumerate}[label=(\Alph*),leftmargin=*,itemsep=0.2em]
\item Listening to the radio
\item Watching a contest
\item Repairing a car battery
\item Attending a conference
\end{enumerate}
\par
\vspace{0.25\baselineskip}
\begin{exampleblock}
\textbf{Transkrip rekonstruksi AI (draft; bukan resmi)}\par
Man: I'd like to rent a compact car for the weekend.\par
Woman: Certainly. May I see your driver's license and credit card?\par
Narrator: Where does this conversation probably take place?\par
\end{exampleblock}
\noindent\textbf{26.}
\begin{enumerate}[label=(\Alph*),leftmargin=*,itemsep=0.2em]
\item An insurance company
\item A car rental agency
\item A real estate agency
\item An apartment complex
\end{enumerate}
\par
\vspace{0.25\baselineskip}
\begin{exampleblock}
\textbf{Transkrip rekonstruksi AI (draft; bukan resmi)}\par
Man: I feel sick, but my exam is today. Do you think I should stay home or go take it?\par
Woman: I wish I could tell you, but I really don't know what you should do.\par
Narrator: What does the woman mean?\par
\end{exampleblock}
\noindent\textbf{27.}
\begin{enumerate}[label=(\Alph*),leftmargin=*,itemsep=0.2em]
\item The man shouldn't attend the exam.
\item The instructor isn't proctoring the exam.
\item In this situation, the man should stay home.
\item She doesn't know what the man should do.
\end{enumerate}
\par
\vspace{0.25\baselineskip}
\begin{exampleblock}
\textbf{Transkrip rekonstruksi AI (draft; bukan resmi)}\par
Man: This course is impossible. I blame the class for my poor grade.\par
Woman: I don't think the course is at fault. It's difficult, but it is fair.\par
Narrator: What does the woman mean?\par
\end{exampleblock}
\noindent\textbf{28.}
\begin{enumerate}[label=(\Alph*),leftmargin=*,itemsep=0.2em]
\item She is a demanding person.
\item She is sympathetic to the man.
\item She doesn't fault the course.
\item She got out of the course.
\end{enumerate}
\par
\vspace{0.25\baselineskip}
\begin{exampleblock}
\textbf{Transkrip rekonstruksi AI (draft; bukan resmi)}\par
Man: Do you have the new edition of the biology textbook?\par
Woman: It's on the shelf with the other required course books. The used copies are in the next aisle.\par
Narrator: Where does this conversation probably take place?\par
\end{exampleblock}
\noindent\textbf{29.}
\begin{enumerate}[label=(\Alph*),leftmargin=*,itemsep=0.2em]
\item A bookstore
\item A publishing house
\item A craft show
\item An art exhibition
\end{enumerate}
\par
\vspace{0.25\baselineskip}
\begin{exampleblock}
\textbf{Transkrip rekonstruksi AI (draft; bukan resmi)}\par
Woman: Congratulations! I heard you came in first yesterday.\par
Man: Thanks. I still can't believe I won the cross-country ski race.\par
Narrator: What did the man do?\par
\end{exampleblock}
\noindent\textbf{30.}
\begin{enumerate}[label=(\Alph*),leftmargin=*,itemsep=0.2em]
\item He walked ten miles.
\item He came here first.
\item He won a ski competition.
\item He lives across the country.
\end{enumerate}
\par
\vspace{0.25\baselineskip}
\subsection*{Part B}
DIRECTIONS: In this part of the test, you will hear longer conversations. After each conversation, you will hear several questions. The conversations and questions will not be repeated. After you hear a question, read the four possible answers in. your test book and select. The best answer. Then, on your answer sheet, find the number of the question and fill in the space that corresponds to the letter of the answer you have chosen. Remember, you are not allowed to take notes or write in your test book.\par
\begin{exampleblock}
Listen to the following example:\par
You will hear:\par
You will read:\par
(A) He has changed jobs.\par
(B) He has two children.\par
(C) He has two jobs.\par
(D) He is looking for a job.\par
From the conversation you learn that Tom has taken an additional job. The best answer to the question "Why is Tom tired?" is (C), "He has two jobs." Therefore the correct answer is (C).\par
\end{exampleblock}
\begin{exampleblock}
\textbf{Transkrip rekonstruksi AI (draft; bukan resmi)}\par
Woman: The movie is almost ready. I'll make some popcorn before we start.\par
Man: Great idea. It's much better to have it here at home than to buy it at a concession stand.\par
Woman: Exactly. Could you get the large cooking pot from the cabinet? We don't need corn on the cob or any special heater, just a pot, some kernels, and a little oil.\par
Man: Popcorn really is one of the cheapest snacks around, isn't it?\par
Woman: It is. That's one reason people like it so much. But I won't add much butter. Butter tastes good, but it mostly adds extra calories.\par
Man: I've always wondered what actually makes popcorn pop.\par
Woman: Each kernel has a tiny amount of moisture inside. When the kernel is heated, that moisture turns to steam. Pressure builds until the kernel explodes.\par
Man: So what is the oil for? Does it make the corn pop?\par
Woman: Not exactly. The oil helps coat the kernels and transfer heat evenly. Without oil, the kernels could burn before they pop.\par
Man: And after they pop, they are huge compared with the original kernels.\par
Woman: Right. A popped kernel is many times larger than it was before.\par
\end{exampleblock}
\noindent\textbf{31.}
\begin{enumerate}[label=(\Alph*),leftmargin=*,itemsep=0.2em]
\item In a movie theater
\item In a baseball stadium
\item At the man and woman's house
\item At a concession stand
\end{enumerate}
\par
\vspace{0.25\baselineskip}
\noindent\textbf{32.}
\begin{enumerate}[label=(\Alph*),leftmargin=*,itemsep=0.2em]
\item Corn on the cob
\item Four types of oil
\item A cooking pot
\item A small heater
\end{enumerate}
\par
\vspace{0.25\baselineskip}
\noindent\textbf{33.}
\begin{enumerate}[label=(\Alph*),leftmargin=*,itemsep=0.2em]
\item They are not good for you.
\item They are cheap.
\item They were invented by settlers.
\item They need moisture.
\end{enumerate}
\par
\vspace{0.25\baselineskip}
\noindent\textbf{34.}
\begin{enumerate}[label=(\Alph*),leftmargin=*,itemsep=0.2em]
\item It's necessary.
\item It adds calories.
\item It's good for you.
\item It tastes terrible.
\end{enumerate}
\par
\vspace{0.25\baselineskip}
\noindent\textbf{35.}
\begin{enumerate}[label=(\Alph*),leftmargin=*,itemsep=0.2em]
\item The liquid coats the corn.
\item Oil puffs up the corn.
\item It makes the kernel explode.
\item Moisture expands at the seams.
\end{enumerate}
\par
\vspace{0.25\baselineskip}
\noindent\textbf{36.}
\begin{enumerate}[label=(\Alph*),leftmargin=*,itemsep=0.2em]
\item You could eat the corn all day long.
\item The corn kernels would burn.
\item The kernels will expand from heating.
\item The corn would have more kernels.
\end{enumerate}
\par
\vspace{0.25\baselineskip}
\noindent\textbf{37.}
\begin{enumerate}[label=(\Alph*),leftmargin=*,itemsep=0.2em]
\item It's many times larger.
\item It's comparatively small.
\item It's hard to measure.
\item It grows steadily.
\end{enumerate}
\par
\vspace{0.25\baselineskip}
\subsection*{Part C}
DIRECTIONS: In Part C you will hear short lectures and extended conversations. At the end of each, you will be asked several questions. Each lecture or conversation and each question will be spoken only one time. For this reason, you must listen carefully to understand what each speaker says. After you hear a question, read the four choices and select the one that best answers the question the speaker asked. Then, on your answer sheet, find the number of the question and blacken the space that contains the letter for the answer you have chosen. Answer all questions according to what is stated or implied in the lecture or conversation.\par
\begin{exampleblock}
Listen to this sample talk.\par
You will hear:\par
Now listen, to the following example.\par
You will hear:\par
You will read:\par
(A) By cars and carriages\par
(B) By bicycles, trains, and carriages\par
(C) On foot and by boat\par
(D) On board ships and trains\par
The best answer to the question "According to the speaker, how did people travel before the invention of the automobile?" is (B), "By bicycles, trains, and carriages." Therefore, the correct answer is (B). You are not allowed to make notes during the test.\par
\end{exampleblock}
\begin{exampleblock}
\textbf{Transkrip rekonstruksi AI (draft; bukan resmi)}\par
Speaker: Welcome to the course. Before we begin the first lecture, I want to review several course policies. Each week's reading and written work must be completed during the week in which it is assigned. Please do not save assignments until the end of the term. If you are absent for more than a few days because of illness or another serious matter, send a notice to my office or to the department. If you miss a test or assignment deadline, you must inform me of the circumstances as soon as possible. Makeup work may be arranged only when I understand the reason for the absence. Otherwise, late assignments will receive lower grades.\par
\end{exampleblock}
\noindent\textbf{38.}
\begin{enumerate}[label=(\Alph*),leftmargin=*,itemsep=0.2em]
\item Course assignments
\item Course policies
\item Class participation
\item Writing projects
\end{enumerate}
\par
\vspace{0.25\baselineskip}
\noindent\textbf{39.}
\begin{enumerate}[label=(\Alph*),leftmargin=*,itemsep=0.2em]
\item By the time the paper is due
\item Over a three-day period
\item Before the end of the term
\item During the week assigned
\end{enumerate}
\par
\vspace{0.25\baselineskip}
\noindent\textbf{40.}
\begin{enumerate}[label=(\Alph*),leftmargin=*,itemsep=0.2em]
\item Assign more homework
\item Take class attendance
\item Give a makeup test
\item Send in a notice
\end{enumerate}
\par
\vspace{0.25\baselineskip}
\noindent\textbf{41.}
\begin{enumerate}[label=(\Alph*),leftmargin=*,itemsep=0.2em]
\item By informing the speaker of the circumstances
\item By talking to their advisers about the class
\item By handing in assignments
\item By taking makeup tests they have missed
\end{enumerate}
\par
\vspace{0.25\baselineskip}
\noindent\textbf{42.}
\begin{enumerate}[label=(\Alph*),leftmargin=*,itemsep=0.2em]
\item Special circumstances may arise at any time.
\item No reasonable excuse will ever be accepted.
\item Assignment grades will be lowered otherwise.
\item Urgent matters must be taken care of in any event.
\end{enumerate}
\par
\vspace{0.25\baselineskip}
\begin{exampleblock}
\textbf{Transkrip rekonstruksi AI (draft; bukan resmi)}\par
Speaker: Today we will discuss how diseases spread from one person to another. A common example is the flu. When a person with the flu coughs, sneezes, or even breathes, tiny particles carrying viruses may be released into the air. Other people can become infected when they breathe in those particles. Bacterial infections can also spread in similar ways, although some diseases may require direct contact or contaminated surfaces. One important point is that people may transmit a virus even before they realize they are ill. For that reason, all infected individuals can potentially spread the disease, not only those who already show clear symptoms.\par
\end{exampleblock}
\noindent\textbf{43.}
\begin{enumerate}[label=(\Alph*),leftmargin=*,itemsep=0.2em]
\item The spread of disease
\item The flu epidemic
\item Flu symptoms
\item Disease variations
\end{enumerate}
\par
\vspace{0.25\baselineskip}
\noindent\textbf{44.}
\begin{enumerate}[label=(\Alph*),leftmargin=*,itemsep=0.2em]
\item The multiplying of viruses
\item Viruses transported by air
\item Drops of contaminated water
\item The regular breathing of air
\end{enumerate}
\par
\vspace{0.25\baselineskip}
\noindent\textbf{45.}
\begin{enumerate}[label=(\Alph*),leftmargin=*,itemsep=0.2em]
\item By adhering to droplets in the air
\item Only by tactile contact
\item Similarly to viral infections
\item Frequently through complications
\end{enumerate}
\par
\vspace{0.25\baselineskip}
\noindent\textbf{46.}
\begin{enumerate}[label=(\Alph*),leftmargin=*,itemsep=0.2em]
\item People who have overt flu symptoms
\item Humans who seek contact
\item Persons who are not self-aware
\item All those infected with a virus
\end{enumerate}
\par
\vspace{0.25\baselineskip}
\begin{exampleblock}
\textbf{Transkrip rekonstruksi AI (draft; bukan resmi)}\par
Speaker: Opinion researchers often study why people choose certain goods and services. In recent surveys, many consumers have said that time is one of their most valuable resources. This helps explain the popularity of fast-food restaurants, household services, and disposable goods, since these products can save time. Researchers also find that appearance, or how people look, strongly affects many buying decisions, especially in areas such as clothing, cosmetics, and exercise programs. But not every consumer choice is practical or appearance-related. Many people spend money on hobbies such as gardening, decorating, cooking, or crafts because these activities allow them to express personal creativity.\par
\end{exampleblock}
\noindent\textbf{47.}
\begin{enumerate}[label=(\Alph*),leftmargin=*,itemsep=0.2em]
\item Product manufacturers
\item Opinion researchers
\item Restaurant owners
\item Commodity traders
\end{enumerate}
\par
\vspace{0.25\baselineskip}
\noindent\textbf{48.}
\begin{enumerate}[label=(\Alph*),leftmargin=*,itemsep=0.2em]
\item Time
\item Money
\item Fast-food restaurants
\item Household services
\end{enumerate}
\par
\vspace{0.25\baselineskip}
\noindent\textbf{49.}
\begin{enumerate}[label=(\Alph*),leftmargin=*,itemsep=0.2em]
\item Lawns
\item Fast-drying paints
\item Looks
\item Disposable goods
\end{enumerate}
\par
\vspace{0.25\baselineskip}
\noindent\textbf{50.}
\begin{enumerate}[label=(\Alph*),leftmargin=*,itemsep=0.2em]
\item To stay healthy
\item To improve their appearance
\item To become energetic
\item To satisfy personal creativity
\end{enumerate}
\par
\vspace{0.25\baselineskip}
\section*{Structure and Written Expression}
\subsection*{Sentence Completion}
This section is designed to test your ability to recognize language structures that are appropriate in standard written English. The questions in this section belong to two types, each of which has special directions.\par
DIRECTIONS: Questions 1-15 are partial sentences. Below each sentence you will see four words or phrases, marked (A), (B), (C), and (D). Select the one word or phrase that best completes the sentence. Then, on your answer sheet, find the number of the question and fill in the space that contains the letter for the answer you have chosen. Fill in the space completely.\par
\begin{exampleblock}
\textbf{EXAMPLE I}\par
\noindent Drying flowers is the best way......them.
\begin{enumerate}[label=(\Alph*),leftmargin=*,itemsep=0.2em]
\item to preserve
\item by preserving
\item preserve
\item preserved
\end{enumerate}
\vspace{0.25\baselineskip}
The sentence should state, "Drying flowers is the best way to preserve them." Therefore, the correct answer is (A).\par
\end{exampleblock}
\begin{exampleblock}
\textbf{EXAMPLE II}\par
\noindent Many American universities......as small, private colleges.
\begin{enumerate}[label=(\Alph*),leftmargin=*,itemsep=0.2em]
\item begun
\item beginning
\item began
\item for the beginning
\end{enumerate}
\vspace{0.25\baselineskip}
The sentence should state, "Many American universities began as small, private colleges." Therefore, the correct answer is (C). After you read the directions, begin work on the questions.\par
\end{exampleblock}
\noindent\textbf{1.} Tennessee has about 140 newspapers,......25 are issued daily.
\begin{enumerate}[label=(\Alph*),leftmargin=*,itemsep=0.2em]
\item about which
\item of which about
\item which are about
\item which is about
\end{enumerate}
\par
\vspace{0.25\baselineskip}
\noindent\textbf{2.} When consumers cannot have everything they want, they have to choose......most.
\begin{enumerate}[label=(\Alph*),leftmargin=*,itemsep=0.2em]
\item they want what
\item what they want
\item they want it
\item that they want
\end{enumerate}
\par
\vspace{0.25\baselineskip}
\noindent\textbf{3.} The temperature of an object rises when......into it.
\begin{enumerate}[label=(\Alph*),leftmargin=*,itemsep=0.2em]
\item heat flow
\item flows hot
\item heat flows
\item hot flow
\end{enumerate}
\par
\vspace{0.25\baselineskip}
\noindent\textbf{4.} From ancient times, people......their land, air and water.
\begin{enumerate}[label=(\Alph*),leftmargin=*,itemsep=0.2em]
\item always have polluting
\item always have pollution
\item have always polluted
\item pollution always has
\end{enumerate}
\par
\vspace{0.25\baselineskip}
\noindent\textbf{5.} Jean Fragonard, was a French artist......portraits of children.
\begin{enumerate}[label=(\Alph*),leftmargin=*,itemsep=0.2em]
\item whose paintings
\item who has painted
\item who painted
\item whose painted
\end{enumerate}
\par
\vspace{0.25\baselineskip}
\noindent\textbf{6.} Overharvesting brought North American alligators, to......in their natural
\begin{enumerate}[label=(\Alph*),leftmargin=*,itemsep=0.2em]
\item nearly extinct
\item near extinction
\item extinct near
\item extinction nearly
\end{enumerate}
\par
\vspace{0.25\baselineskip}
\noindent\textbf{7.} The Ford Foundation was established in 1936 to advance human well-being by......funds for education.
\begin{enumerate}[label=(\Alph*),leftmargin=*,itemsep=0.2em]
\item contribute
\item contribution
\item to contribute
\item contributing
\end{enumerate}
\par
\vspace{0.25\baselineskip}
\noindent\textbf{8.} Fireweed received its name because it......after a forest fire.
\begin{enumerate}[label=(\Alph*),leftmargin=*,itemsep=0.2em]
\item quick growth
\item grows quickly
\item quickly grown
\item growing quickly
\end{enumerate}
\par
\vspace{0.25\baselineskip}
\noindent\textbf{9.} Florida's long coastline and warm weather......swimmers to its sandy shores.
\begin{enumerate}[label=(\Alph*),leftmargin=*,itemsep=0.2em]
\item attracts
\item attract
\item they attract
\item is attracted by
\end{enumerate}
\par
\vspace{0.25\baselineskip}
\noindent\textbf{10.} Amazon pygmies consider their songs......part of their culture.
\begin{enumerate}[label=(\Alph*),leftmargin=*,itemsep=0.2em]
\item an important and extremely
\item as extreme and important
\item an extremely important
\item as extreme importance
\end{enumerate}
\par
\vspace{0.25\baselineskip}
\noindent\textbf{11.} Psychologists define anxiety as a feeling of dread, apprehension, or......
\begin{enumerate}[label=(\Alph*),leftmargin=*,itemsep=0.2em]
\item afraid
\item be afraid
\item having fear
\item fear
\end{enumerate}
\par
\vspace{0.25\baselineskip}
\noindent\textbf{12.} Ancient nations have used......on their emblems and flags for thousands of years.
\begin{enumerate}[label=(\Alph*),leftmargin=*,itemsep=0.2em]
\item the same symbols as
\item the same symbols
\item symbols the same as
\item symbols as the same
\end{enumerate}
\par
\vspace{0.25\baselineskip}
\noindent\textbf{13.} For years experts......the effect of coaching and preparatory courses on test scores.
\begin{enumerate}[label=(\Alph*),leftmargin=*,itemsep=0.2em]
\item are examining
\item had been examined
\item have been examining.
\item having been examined
\end{enumerate}
\par
\vspace{0.25\baselineskip}
\noindent\textbf{14.} The nitrogen cycle......of nitrogen through the atmosphere, hydrosphere, and lithosphere.
\begin{enumerate}[label=(\Alph*),leftmargin=*,itemsep=0.2em]
\item the circulation is
\item is the circulation
\item it is the circulation
\item is it the circulation
\end{enumerate}
\par
\vspace{0.25\baselineskip}
\noindent\textbf{15.} While working as a clerk, Edison spent much of his time......the stock ticker.
\begin{enumerate}[label=(\Alph*),leftmargin=*,itemsep=0.2em]
\item he studied
\item to study
\item on study
\item studying
\end{enumerate}
\par
\vspace{0.25\baselineskip}
\subsection*{Error Identification}
DIRECTIONS: In questions 16-40 every sentence has four words or phrases that are underlined. The four underlined portions of each sentence are marked (A), (B), (C), and (D). Identify the one word or phrase that makes the sentence incorrect. Then, on your answer sheet, find the number of the question and fill in the space that contains the letter for the answer you have chosen. Fill in the space completely.\par
\begin{exampleblock}
\textbf{EXAMPLE I}\par
Christopher Columbus \uline{has sailed}(A) from Europe in 1492 and \uline{discovered a}(B) \uline{new land}(C) he \uline{thought to}(D) be India.\par
The sentence should state, "Christopher Columbus sailed from Europe in 1492 and discovered a new land he thought to be India." Therefore, you should choose answer (A).\par
\end{exampleblock}
\begin{exampleblock}
\textbf{EXAMPLE II}\par
\uline{As the roles}(A) of people in society change, \uline{so does}(B) the \uline{rules of}(C) conduct in certain \uline{situations}(D).\par
The sentence should state, "As the roles of people in society change, so do the rules of\par
conduct in certain situations." Therefore, you should choose answer (B).\par
\end{exampleblock}
\noindent\textbf{16.} James Maxwell \uline{based}(A) his \uline{work on}(B) the \uline{discoveries}(C) of the English \uline{physical}(D) Michael Faraday.
\par
\vspace{\questionblockgap}
\noindent\textbf{17.} The behavior of animals \uline{appears to}(A) depend on \uline{patterns of}(B) \uline{reactions which}(C) \uline{they are}(D) born.
\par
\vspace{\questionblockgap}
\noindent\textbf{18.} \uline{For centuries}(A), people \uline{have wondered}(B) why \uline{have they}(C) particular dreams while \uline{they sleep}(D).
\par
\vspace{\questionblockgap}
\noindent\textbf{19.} \uline{After the}(A) Roman Empire \uline{has collapsed}(B), Europe had \uline{no regular}(C) \uline{postal}(D) service.
\par
\vspace{\questionblockgap}
\noindent\textbf{20.} Most frequently, asphalt \uline{is used}(A) \uline{to pave}(B) streets, \uline{avenues}(C), highways, and \uline{airport}(D).
\par
\vspace{\questionblockgap}
\noindent\textbf{21.} Samuel Coleridge was \uline{poet}(A) and philosopher of \uline{the English}(B) \uline{romantic}(C) movement in \uline{the early}(D) 1800s.
\par
\vspace{\questionblockgap}
\noindent\textbf{22.} When one person \uline{intentional}(A) takes the property of \uline{another}(B) without legal \uline{justification}(C), the crime is \uline{called}(D) theft.
\par
\vspace{\questionblockgap}
\noindent\textbf{23.} \uline{Computers}(A) programs, \uline{catalogues}(B), directories, and \uline{collections}(C) of data are \uline{protected}(D) by copyright laws.
\par
\vspace{\questionblockgap}
\noindent\textbf{24.} The water temperature in \uline{a spring}(A) depends \uline{on that}(B) of the soil through \uline{where}(C) the water \uline{flows}(D).
\par
\vspace{\questionblockgap}
\noindent\textbf{25.} Tree squirrels are active, \uline{noisy}(A), and \uline{lively}(B) \uline{animals that}(C) make \uline{its home}(D) in tree trunks.
\par
\vspace{\questionblockgap}
\noindent\textbf{26.} Botanists \uline{have determined}(A) that \uline{there is}(B) more than 60 \uline{species}(C) of \uline{sunflowers}(D).
\par
\vspace{\questionblockgap}
\noindent\textbf{27.} Pure sodium \uline{immediately}(A) \uline{combines}(B) with oxygen when \uline{is exposed}(C) \uline{to air}(D).
\par
\vspace{\questionblockgap}
\noindent\textbf{28.} Special education is intended \uline{help}(A) \uline{both}(B) handicapped and gifted children to \uline{reach}(C) their \uline{learning}(D) potentials.
\par
\vspace{\questionblockgap}
\noindent\textbf{29.} \uline{Trading}(A) fairs \uline{held}(B) in Antwerp during the 1300s \uline{brought}(C) \uline{famous}(D) to the city.
\par
\vspace{\questionblockgap}
\noindent\textbf{30.} By 1850, California had \uline{a sufficiently}(A) \uline{large population}(B) \uline{to recognized}(C) \uline{as a state}(D).
\par
\vspace{\questionblockgap}
\noindent\textbf{31.} Gingham is \uline{a fabric}(A) used to make \uline{dresses}(B), \uline{curtains}(C), and \uline{furnitures}(D) covers.
\par
\vspace{\questionblockgap}
\noindent\textbf{32.} Plato's most \uline{last}(A) contribution to mathematics was his \uline{insistence}(B) on \uline{using}(C) \uline{reasoning}(D) in geometry.
\par
\vspace{\questionblockgap}
\noindent\textbf{33.} Sedatives are \uline{a group}(A) of drugs \uline{that legally}(B) prescribed and should \uline{be taken}(C) \uline{as directed}(D).
\par
\vspace{\questionblockgap}
\noindent\textbf{34.} Maps \uline{that show}(A) \uline{detail}(B) landforms \uline{are commonly}(C) \uline{used}(D) in physical geography and geology.
\par
\vspace{\questionblockgap}
\noindent\textbf{35.} When the minerals \uline{needed}(A) for corn \uline{to grow}(B) \uline{are lack}(C), the husks may \uline{be stunted}(D).
\par
\vspace{\questionblockgap}
\noindent\textbf{36.} Rembrandt was \uline{forced}(A) to \uline{declare}(B) bankruptcy in 1656, and his \uline{possessions}(C) were \uline{sale}(D).
\par
\vspace{\questionblockgap}
\noindent\textbf{37.} Pearls and \uline{similar}(A) substances \uline{may be}(B) \uline{classified}(C) by \uline{how are}(D) cultivated.
\par
\vspace{\questionblockgap}
\noindent\textbf{38.} During \uline{winter}(A), grizzly bears \uline{live in}(B) dens, caves or \uline{the other}(C) natural \uline{shelters}(D).
\par
\vspace{\questionblockgap}
\noindent\textbf{39.} The neck of a \uline{classical}(A) guitar is \uline{wider}(B) than \uline{those}(C) of a steel-\uline{string}(D) guitar.
\par
\vspace{\questionblockgap}
\noindent\textbf{40.} \uline{Insufficient}(A) oxygen causes lactic acid to \uline{built up}(B) in the \uline{muscles}(C) of long-\uline{distance}(D) runners.
\par
\vspace{\questionblockgap}
\section*{Reading Comprehension}
\begin{center}
\textbf{Time 55 minutes}
\end{center}
DIRECTIONS: In this section you will read several passages. Each passage is followed by a series of questions. For questions 1-50, you need to select the best answer, (A), (B), (C), or (D), to each question. Then, on your answer sheet, find the number of the question and fill in the space that contains the letter of the answer you have selected. Fill in the space completely. Answer all questions following a passage on the basis of what is stated or implied in the passage.\par
\begin{exampleblock}
\small
\sloppy
Read the following passage:\par
A tomahawk is a small ax used as a tool and a weapon by the North American Indian\par
tribes. An average tomahawk was not very long and did not weigh a great deal. Originally,\par
the head of the tomahawk was made of a shaped stone or an animal bone and was mounted on\par
Line a wooden handle. After the arrival of the European settlers, the Indians began to use toma-\par
hawks with iron heads. Indian males and females of all ages used tomahawks to chop and cut wood, pound stakes into the ground to put up wigwams, and do many other chores. Indian warriors relied on tomahawks as weapons and even threw them at their enemies. Some types of tomahawks were used in religious ceremonies. Contemporary American idioms reflect\par
this aspect of American heritage.\par
\end{exampleblock}
\begin{exampleblock}
\textbf{EXAMPLE I}\par
\noindent Early tomahawk heads were made of
\begin{enumerate}[label=(\Alph*),leftmargin=*,itemsep=0.2em]
\item stone or bone
\item wood or sticks
\item European iron
\item religious weapons
\end{enumerate}
\vspace{0.25\baselineskip}
According to the passage, early tomahawk heads were made of stone or bone. Therefore, the correct answer is (A).\par
\end{exampleblock}
\begin{exampleblock}
\textbf{EXAMPLE II}\par
\noindent How has the Indian use of tomahawks affected American daily life today?
\begin{enumerate}[label=(\Alph*),leftmargin=*,itemsep=0.2em]
\item Tomahawks are still used as weapons.
\item Tomahawks are used as tools for.certain jobs.
\item Contemporary language refers to tomahawks.
\item Indian tribes cherish tomahawks as heirlooms.
\end{enumerate}
\vspace{0.25\baselineskip}
The passage states, "Contemporary American idioms reflect this aspect of American heritage." The correct answer is (C). After you read the directions, begin work on the questions.\par
\end{exampleblock}
\subsection*{Questions 1-10}
\begingroup
\small
\sloppy
\setlength{\parskip}{0pt}
\setlength{\parindent}{0pt}
During the Middle Ages, societies were based on military relationships, as landowners\par
formed their own foot armies into which they drafted their tenants and hired hands. The in-\par
fantry that fought its way forward against the opposition engaged in heavy ground battles\par
that proved costly in the ratio of losses to wins. These soldiers carried darts, javelins, and\par
slings to be used before closing ranks with the enemy, although their swords and halberds\par
 delivered crushing blows on contact. Such armed forces were active for limited periods of\par
time and had a predominantly defensive function, displayed in hand-to-hand combat.\par
Because this sporadic and untrained organization was ineffective, the ruling classes be-\par
gan to hire mercenaries who were generously compensated for their tasks and subject to con-\par
tractual terms of agreement. The greatest idiosyncrasy of a hired military force was that the\par
troops sometimes deserted their employers if they could bank on a higher remuneration\par
from the opposition. The Swiss pikemen became the best-known mercenaries of the late\par
Middle Ages. In the 1300s, they practically invented a crude body armor of leather and\par
quilted layered head gear with nose and skull plates, ornamented with crests. Their tower\par
shields proved indispensable against a shower of arrows, and their helmets progressed from\par
cone cups to visors hinged at the temples. As their notoriety increased, so did their wages,\par
and eventually they were rounded into military companies that later grew into the basic\par
units in almost all armies. During the same period, the first full-size army of professional\par
soldiers emerged in the Ottoman Empire. What set these troops apart from other contempo-\par
rary armies was that these soldiers remained on duty in peacetime.\par
Companies of mercenaries were employed on a permanent basis in 1445, when King\par
Charles VII created a regular military organization, complete with a designated hierarchy.\par
Gunpowder accelerated the emergence of military tactics and strategy that ultimately af-\par
fected the conceptualization of war on a broad scale. Cannons further widened the gap be-\par
tween the attacking and the defending lineups, and undermined the exclusivity of contact\par
battles.\par
\endgroup
\medskip
\noindent\textbf{1.} What is the main purpose of the passage?
\begin{enumerate}[label=(\Alph*),leftmargin=*,itemsep=0.2em]
\item To distinguish between laborers and mercenaries
\item To change the existing view of the military
\item To cite examples of armor in the Middle Ages
\item To trace the origins of military organization
\end{enumerate}
\par
\vspace{0.25\baselineskip}
\noindent\textbf{2.} In line 4, the word "ratio" is closest in meaning to
\begin{enumerate}[label=(\Alph*),leftmargin=*,itemsep=0.2em]
\item quota
\item reason
\item proportion
\item pace
\end{enumerate}
\par
\vspace{0.25\baselineskip}
\noindent\textbf{3.} Which of the following statements can be inferred from the first paragraph?
\begin{enumerate}[label=(\Alph*),leftmargin=*,itemsep=0.2em]
\item Temporary armies of farmers were not well trained.
\item Drafting farmers into armies was costly.
\item Heavy ground battles were won during combat.
\item Infantry was directed to the opposition for support.
\end{enumerate}
\par
\vspace{0.25\baselineskip}
\noindent\textbf{4.} In line 8, the word "sporadic" is closest in meaning to
\begin{enumerate}[label=(\Alph*),leftmargin=*,itemsep=0.2em]
\item spirited
\item splendid
\item irreverent
\item irregular
\end{enumerate}
\par
\vspace{0.25\baselineskip}
\noindent\textbf{5.} Which of the following statements about the Swiss pikemen is supported by the passage?
\begin{enumerate}[label=(\Alph*),leftmargin=*,itemsep=0.2em]
\item Their weapons and skills were ahead of their time.
\item Their gear ensured their fame as well-dressed soldiers.
\item The demand for their cavalry made them the best-paid army.
\item Their weapons were issued to nonprofessionals as well.
\end{enumerate}
\par
\vspace{0.25\baselineskip}
\noindent\textbf{6.} Where in the passage does the author state the reasons for the emergence of professional armies?
\begin{enumerate}[label=(\Alph*),leftmargin=*,itemsep=0.2em]
\item Lines 1-4
\item Lines 8-10
\item Lines 12-14
\item Lines 15-18
\end{enumerate}
\par
\vspace{0.25\baselineskip}
\noindent\textbf{7.} The author of the passage implies that the soldiers in mercenary armies were
\begin{enumerate}[label=(\Alph*),leftmargin=*,itemsep=0.2em]
\item not loyal
\item not effective
\item well guarded
\item well rounded
\end{enumerate}
\par
\vspace{0.25\baselineskip}
\noindent\textbf{8.} In line 19, the phrase "these troops" refers to
\begin{enumerate}[label=(\Alph*),leftmargin=*,itemsep=0.2em]
\item the Swiss pikemen
\item military companies
\item almost all armies
\item Ottoman soldiers
\end{enumerate}
\par
\vspace{0.25\baselineskip}
\noindent\textbf{9.} According to the passage, the first army of professionals was mobilized
\begin{enumerate}[label=(\Alph*),leftmargin=*,itemsep=0.2em]
\item only in peacetime
\item only in wartime
\item in times of anticipated war
\item both during war and during peace
\end{enumerate}
\par
\vspace{0.25\baselineskip}
\noindent\textbf{10.} In line 25, the word "undermined" is closest in meaning to
\begin{enumerate}[label=(\Alph*),leftmargin=*,itemsep=0.2em]
\item underestimated
\item reduced
\item undersized
\item shred
\end{enumerate}
\par
\vspace{0.25\baselineskip}
\subsection*{Questions 11-18}
\begingroup
\small
\sloppy
\setlength{\parskip}{0pt}
\setlength{\parindent}{0pt}
Observations of nature gained a foothold in art in the 1860s and 1870s when painters in-\par
terested in science attempted to analyze the effects of light on color by means of physics. If\par
the goal of impressionist painters was to copy the visual qualities of sunlight at different an-\par
gles, they needed to reproduce light as it appears to the spectator when reflected from the\par
surfaces of structures. In painting, the effects of shade were conveyed by using small strokes\par
to minimize breaks between hues. The so-called divided color method appeared to grasp a\par
shimmering reflection of shadows when minimal portions of primary-color paints were ap-\par
plied directly to the canvas, instead of being blended on the palette.\par
Edouard Manet departed from the fairy-tale style of painting with its tacit symbolism\par
and centered his compositions around the visual reality of ordinary objects. Mary Cassatt fol-\par
lowed with her spontaneous and subtle portraits of children, and Edgar Degas depicted bal-\par
let dancers in their artful poses and the color schemes of their costumes in soft colors.\par
Postimpressionism built on the techniques developed by impressionists and supple-\par
mented it with keen insight into other dimensions of objects and scenes. Paul Gaugin chose\par
to disregard the classical conventions of composition, the application of color, and the shap-\par
ing of form and imitated primitivist art that upheld the beauty of native drawings in Tahiti.\par
Henri Matisse, created a unique style of poster graphics, deceptively simplistic in its rhythm\par
and texture. In his view, paintings were intended to brighten and improve reality, not copy\par
it. He noted that photography can accomplish this latter goal just as well, or even better.\par
\endgroup
\medskip
\noindent\textbf{11.} This passage probably comes from a longer work on
\begin{enumerate}[label=(\Alph*),leftmargin=*,itemsep=0.2em]
\item science and the fine arts
\item great masters of impressionist painting
\item light in the paintings of the 1800s
\item new techniques in the art of the 1800s
\end{enumerate}
\par
\vspace{0.25\baselineskip}
\noindent\textbf{12.} The author of the passage implies that the goal of the impressionists was to
\begin{enumerate}[label=(\Alph*),leftmargin=*,itemsep=0.2em]
\item reproduce light and shadow exactly as they appeared
\item depict light as it appeared on different surfaces
\item copy the pattern of sun rays at different angles
\item reflect light and shadow from surfaces of structures
\end{enumerate}
\par
\vspace{0.25\baselineskip}
\noindent\textbf{13.} In line 7, the word "shimmering" is closest in meaning to
\begin{enumerate}[label=(\Alph*),leftmargin=*,itemsep=0.2em]
\item gleaming
\item strong
\item trembling
\item spotless
\end{enumerate}
\par
\vspace{0.25\baselineskip}
\noindent\textbf{14.} In line 8, the word "blended" is closest in meaning to
\begin{enumerate}[label=(\Alph*),leftmargin=*,itemsep=0.2em]
\item replenished
\item placed
\item rendered
\item mixed
\end{enumerate}
\par
\vspace{0.25\baselineskip}
\noindent\textbf{15.} What technique did painters employ to represent light as it appeared to the artist?
\begin{enumerate}[label=(\Alph*),leftmargin=*,itemsep=0.2em]
\item They used as little paint as was necessary.
\item They graded the shades of color in hues.
\item Their paint colors were dark and muted.
\item Their brush strokes were slow and cautious.
\end{enumerate}
\par
\vspace{0.25\baselineskip}
\noindent\textbf{16.} It is implied in the passage that in the late 1800s artists painted
\begin{enumerate}[label=(\Alph*),leftmargin=*,itemsep=0.2em]
\item fantastic images and swirling action
\item luxurious ornaments and shapely figures
\item unusual textures and luscious colors
\item ordinary objects in varied intensities of light
\end{enumerate}
\par
\vspace{0.25\baselineskip}
\noindent\textbf{17.} It can be inferred from the passage that postimpressionism
\begin{enumerate}[label=(\Alph*),leftmargin=*,itemsep=0.2em]
\item developed into a trend parallel to impressionism
\item followed impressionism in the development of technique
\item flourished independently of impressionism
\item negated impressionist insight into objects and scenes
\end{enumerate}
\par
\vspace{0.25\baselineskip}
\noindent\textbf{18.} According to the passage, Matisse saw the purpose of his art as
\begin{enumerate}[label=(\Alph*),leftmargin=*,itemsep=0.2em]
\item depicting life realistically
\item enhancing life and reality
\item showing life through photography
\item imparting rhythm to drawings
\end{enumerate}
\par
\vspace{0.25\baselineskip}
\subsection*{Questions 19-30}
\begingroup
\small
\sloppy
\setlength{\parskip}{0pt}
\setlength{\parindent}{0pt}
The symptoms of hay fever include watery and itchy eyes and a runny, congested nose.\par
People suffering from hay fever may experience occasional wheezing and repeated bouts of\par
sneezing and may even lose their sense of smell. Some victims of hay fever may also have\par
stopped-up ears. About 30 percent of those who suffer from hay fever may develop the symp-\par
toms associated with periodic asthma or a sinus infection. The allergen-antibody theory\par
does not fully explain allergic reactions because the membranes and glands in eyes and ears\par
are controlled by the independent nervous system, which keeps these organs in balance. But\par
the independent nervous system itself is part of the emotional-response center and may\par
cause the feelings of anger, fear, resentment, and lack of self-confidence in reaction to al-\par
lergy-causing substances.\par
The most common cause of hay fever is the pollen of ragweed, which blossoms during the\par
summer and autumn. When airborne pollen particles, as well as mold, come into contact\par
with the victim's membranes, they can cause allergic reactions that release histamine and re-\par
sult in a virtual blockage of air passages. To prevent hay fever or to decrease the severity of its\par
symptoms, contact with the ragweed pollen should be reduced. Although some communi-\par
ties have attempted to eliminate the plants that cause the reactions, elimination programs\par
have not been successful because airborne pollen can travel considerable distances. Antihis-\par
tamine can help with short but severe attacks. Over extended periods of time, however, pa-\par
tients are prescribed a series of injections of the substance to which they are sensitive in order\par
to increase immunity and thus be relieved of the seasonal allergy\par
\endgroup
\medskip
\noindent\textbf{19.} It can be inferred from the passage that the phrase "hay fever" refers to
\begin{enumerate}[label=(\Alph*),leftmargin=*,itemsep=0.2em]
\item fodder for cattle
\item a seasonal discomfort
\item viral bacteria
\item a lung disease
\end{enumerate}
\par
\vspace{0.25\baselineskip}
\noindent\textbf{20.} According to the passage, the symptoms of the allergy are predominantly
\begin{enumerate}[label=(\Alph*),leftmargin=*,itemsep=0.2em]
\item abdominal
\item intestinal
\item respiratory
\item chronic
\end{enumerate}
\par
\vspace{0.25\baselineskip}
\noindent\textbf{21.} What can be inferred from the first paragraph?
\begin{enumerate}[label=(\Alph*),leftmargin=*,itemsep=0.2em]
\item Hay fever may cause severe allergic reactions and even death.
\item The cause of allergic reactions has not been determined.
\item The nervous system balances allergic reactions.
\item People should not have an emotional response to allergic reactions.
\end{enumerate}
\par
\vspace{0.25\baselineskip}
\noindent\textbf{22.} According to the passage, patients suffering from hay fever may also experience
\begin{enumerate}[label=(\Alph*),leftmargin=*,itemsep=0.2em]
\item hunger pains
\item mood swings
\item nervous blockages
\item sensory perceptions
\end{enumerate}
\par
\vspace{0.25\baselineskip}
\noindent\textbf{23.} In line 9, the word "resentment" is closest in meaning to
\begin{enumerate}[label=(\Alph*),leftmargin=*,itemsep=0.2em]
\item reprieve
\item reprisal
\item acrimony
\item grief
\end{enumerate}
\par
\vspace{0.25\baselineskip}
\noindent\textbf{24.} It can be inferred from the passage that a frequent source of allergy-causing irritants can be
\begin{enumerate}[label=(\Alph*),leftmargin=*,itemsep=0.2em]
\item organic matter
\item larynx infections
\item human contact
\item ear membranes
\end{enumerate}
\par
\vspace{0.25\baselineskip}
\noindent\textbf{25.} According to the passage, the irritants are transported by
\begin{enumerate}[label=(\Alph*),leftmargin=*,itemsep=0.2em]
\item wind
\item food
\item travelers
\item air passages
\end{enumerate}
\par
\vspace{0.25\baselineskip}
\noindent\textbf{26.} In line 14, the word "blockage" is closest in meaning to
\begin{enumerate}[label=(\Alph*),leftmargin=*,itemsep=0.2em]
\item obstruction
\item bleeding
\item enlargement
\item dryness
\end{enumerate}
\par
\vspace{0.25\baselineskip}
\noindent\textbf{27.} According to the passage, to avoid incidents of hay fever, patients need to
\begin{enumerate}[label=(\Alph*),leftmargin=*,itemsep=0.2em]
\item avoid interactions with other patients
\item avoid exposure to pollen
\item increase their self-confidence
\item take doses of prescribed medicine
\end{enumerate}
\par
\vspace{0.25\baselineskip}
\noindent\textbf{28.} Which of the following is not mentioned in the passage as a cause of allergies?
\begin{enumerate}[label=(\Alph*),leftmargin=*,itemsep=0.2em]
\item pollen
\item mold
\item flowers
\item injections
\end{enumerate}
\par
\vspace{0.25\baselineskip}
\noindent\textbf{29.} It can be inferred from the passage that hay fever
\begin{enumerate}[label=(\Alph*),leftmargin=*,itemsep=0.2em]
\item has no effective antibodies
\item has no known cure
\item is rooted in the human psyche
\item can be likened to a breakdown
\end{enumerate}
\par
\vspace{0.25\baselineskip}
\noindent\textbf{30.} A paragraph following this passage would most probably discuss
\begin{enumerate}[label=(\Alph*),leftmargin=*,itemsep=0.2em]
\item how the nervous system alerts patients
\item how the immune system reacts to allergens
\item what other diseases can be relieved by vaccines
\item what flowers are harmless to hay fever patients
\end{enumerate}
\par
\vspace{0.25\baselineskip}
\subsection*{Questions 31-40}
\begingroup
\small
\sloppy
\setlength{\parskip}{0pt}
\setlength{\parindent}{0pt}
Prehistoric horses were far removed from the horses that Christopher Columbus brought\par
on his ships during his second voyage to the New World. Although fossil remains of "dawn\par
horses" have been excavated in several sites in Wyoming and New Mexico, these animals,\par
which were biologically different from contemporary horses, had become extinct millennia\par
before the onset of the Indian era. Although moviegoers visualize an Indian as a horse rider,\par
Indians were not familiar with horses until the Spanish brought them to Mexico, New Mex-\par
ico, Florida, and the West Indies in 1519. Those that escaped from the conquerors or were\par
left behind became the ancestors of the wild horses that still roam the southwestern regions\par
of the country. The Indian tribes scattered in the western plains began to breed horses about\par
1600.\par
The arrival of the horse produced a ripple effect throughout the Great Plains as the Indi-\par
ans living there were not nomadic and were engaged in rudimentary farming and graze-\par
land hunting. Tracking stampeding herds of buffalo and elk on foot was not the best way to\par
stock quantities of meat to adequately feed the entire tribe during the winter. However,\par
mounted on horses, the hunting teams could cover ground within a substantial distance\par
from their camps and transport their game back to be roasted, dried into jerky, or smoked\par
for preservation. The hunters responsible for tribe provisions stayed on the move almost con-\par
tinuously, replacing their earth-and-sod lodges with tepees. Horses carried not only their\par
riders but also their possessions and booty. The Blackfoot Indians of the Canadian plains\par
turned almost exclusively hunters, and the Crow split off from the mainstream Indian farm-\par
ing in favor of hunting. In fact, some of the Apache splinter groups abandoned agricultural\par
cultivation altogether.\par
The horse also drastically altered Indian warfare by allowing rapid maneuvering before,\par
during, and after skirmishes. With the advent of the horse, the Apache, Arapahoe, and\par
Cheyenne established themselves as a territorial monopoly in the Plains. Because Indians\par
did not have the wheel and had dragged their belongings from one settlement to another,\par
horses also enabled them to become more mobile and expedient during tribal migrations. In\par
fact, the Cheyenne abolished the custom of discarding belongings and tepee skins simply\par
because there were no means to transport them\par
\endgroup
\medskip
\noindent\textbf{31.} In line 3, the word "excavated" is closest in meaning to
\begin{enumerate}[label=(\Alph*),leftmargin=*,itemsep=0.2em]
\item exasperated
\item extinguished
\item hunted down
\item dug up
\end{enumerate}
\par
\vspace{0.25\baselineskip}
\noindent\textbf{32.} According to the passage, how many genetic species of horses are known today?
\begin{enumerate}[label=(\Alph*),leftmargin=*,itemsep=0.2em]
\item One
\item Two
\item Three
\item Four
\end{enumerate}
\par
\vspace{0.25\baselineskip}
\noindent\textbf{33.} In line 7, the word "those" refers to
\begin{enumerate}[label=(\Alph*),leftmargin=*,itemsep=0.2em]
\item West Indies
\item The Spanish
\item horses,
\item Indians
\end{enumerate}
\par
\vspace{0.25\baselineskip}
\noindent\textbf{34.} According to the passage, American Indians
\begin{enumerate}[label=(\Alph*),leftmargin=*,itemsep=0.2em]
\item tamed horses in the early 1500s
\item farmed with horses in the 1500s
\item were exposed to horses in the 1500s
\item have ridden horses since prehistoric times
\end{enumerate}
\par
\vspace{0.25\baselineskip}
\noindent\textbf{35.} The author of the passage probably believes that the popular image of American Indians before the arrival of Europeans
\begin{enumerate}[label=(\Alph*),leftmargin=*,itemsep=0.2em]
\item is not theoretically viable
\item cannot be realistically described
\item cannot be discussed briefly
\item is not historically accurate
\end{enumerate}
\par
\vspace{0.25\baselineskip}
\noindent\textbf{36.} According to the passage, after the arrival of Europeans, the Indian tribes inhabiting the Great Plains
\begin{enumerate}[label=(\Alph*),leftmargin=*,itemsep=0.2em]
\item herded undomesticated buffalo
\item played complicated hunting games
\item had sedentary and tranquil life-styles
\item improved their hunting techniques
\end{enumerate}
\par
\vspace{0.25\baselineskip}
\noindent\textbf{37.} In line 17, the word "provisions" is closest in meaning to
\begin{enumerate}[label=(\Alph*),leftmargin=*,itemsep=0.2em]
\item supplies
\item health
\item weapons
\item attire
\end{enumerate}
\par
\vspace{0.25\baselineskip}
\noindent\textbf{38.} According to the passage, American Indians invented various methods for
\begin{enumerate}[label=(\Alph*),leftmargin=*,itemsep=0.2em]
\item dislocating their traps
\item communicating over great distances
\item conducting their hostile excursions
\item keeping their possessions
\end{enumerate}
\par
\vspace{0.25\baselineskip}
\noindent\textbf{39.} It can be inferred from the passage that Indians did NOT
\begin{enumerate}[label=(\Alph*),leftmargin=*,itemsep=0.2em]
\item accrue tribal wealth
\item assign sustenance tasks
\item pursue stampedes
\item use covered wagons
\end{enumerate}
\par
\vspace{0.25\baselineskip}
\noindent\textbf{40.} It can be inferred from the passage that the arrival of horses in the Americas
\begin{enumerate}[label=(\Alph*),leftmargin=*,itemsep=0.2em]
\item led to the dispersal of Indian tribes throughout the continent
\item made Indian tribes relinquish their territorial monopolies
\item altered the future course of the Indian way of life
\item shattered the advancement of the Indian culture
\end{enumerate}
\par
\vspace{0.25\baselineskip}
\subsection*{Questions 41-50}
\begingroup
\small
\sloppy
\setlength{\parskip}{0pt}
\setlength{\parindent}{0pt}
Lighthouses and lightships employ signal lights and foghorns to warn boaters of shoals\par
and of oncoming foul weather. In recreational marinas, Coast Guard stations and yacht\par
clubs inform boaters of weather and water conditions with storm flag signals displayed dur-\par
ing the daytime. Amateur boaters are required to acquaint themselves with the signals that\par
can make them aware of oncoming storms and small craft advisories. Boaters who can recog-\par
nize an approaching storm when no warnings are posted have the advantage of time when\par
heading for shelter. Compasses and marine charts identify the locations of anchored whistles\par
and floating buoys that have battery-powered lights, bells, or horns to sound or show warn-\par
ings with the onset of high winds. On the marine charts, radio beacons and flashing buoys\par
are assigned numbers that help navigators to identify their locations and to mark the edges\par
of channels, underwater obstructions, reefs, and wrecks.\par
In each section of the coastline, lighthouses emit characteristic lights, published in light\par
lists, so that mariners can zero in on their bearings by observing the beam pattern and con-\par
sulting the list. "Making lights" signal approaching vessels to make land, and "leading\par
lights" guide navigators into bays and harbors along navigable waterways.\par
The lighthouse tower, constructed on solid rock and pneumatic caissons, contains light-\par
ing mechanisms, engines, and spare parts, as well as the keeper's quarters. Saucerlike Fres-\par
nel lenses with pentagonal prisms project light at irregular, alternating intervals, while\par
sealed-beam lenses rotate at varying speeds, similarly to the traditional long-range search\par
light. The older, barrel-shaped lenses and kerosene-burning lamps, made up of prisms and\par
glass panels attached to a metal frame, have been replaced by acetylene gas burners and in-\par
candescent lamps that can be operated either manually or automatically and that shut off at\par
daybreak\par
\endgroup
\medskip
\noindent\textbf{41.} What is the best title for the passage?
\begin{enumerate}[label=(\Alph*),leftmargin=*,itemsep=0.2em]
\item The Function of Naval Signals and Buoys
\item Marine Installations for Boaters' Safety
\item Lighthouses and Maritime Signaling Devices
\item Lighthouses and Occulting Instrumentation
\end{enumerate}
\par
\vspace{0.25\baselineskip}
\noindent\textbf{42.} According to the passage, what is the purpose of lighthouses?
\begin{enumerate}[label=(\Alph*),leftmargin=*,itemsep=0.2em]
\item To measure weather and water conditions
\item To lead navigators to their destination
\item To flag down oncoming vessels
\item To display light-coded messages
\end{enumerate}
\par
\vspace{0.25\baselineskip}
\noindent\textbf{43.} According to the passage, buoys can be best described as
\begin{enumerate}[label=(\Alph*),leftmargin=*,itemsep=0.2em]
\item hanging constructions
\item electrical horns
\item floating devices
\item underwater anchors
\end{enumerate}
\par
\vspace{0.25\baselineskip}
\noindent\textbf{44.} Why do recreational mariners need to be familiar with the Coast Guard signaling system?
\begin{enumerate}[label=(\Alph*),leftmargin=*,itemsep=0.2em]
\item To learn navigation and rowing
\item To chart the course of storms
\item To learn about oncoming storms
\item To overtake rivals in boating races
\end{enumerate}
\par
\vspace{0.25\baselineskip}
\noindent\textbf{45.} In line 6, the phrase "the advantage of time" means most nearly
\begin{enumerate}[label=(\Alph*),leftmargin=*,itemsep=0.2em]
\item an early warning
\item a stable bearing
\item a timed landing
\item a timely arrival
\end{enumerate}
\par
\vspace{0.25\baselineskip}
\noindent\textbf{46.} The author of the passage implies that flashing lights from lighthouses function as
\begin{enumerate}[label=(\Alph*),leftmargin=*,itemsep=0.2em]
\item a characteristic marine elevation
\item beams with a constant rotating speed
\item an indication of the lens thickness
\item an identification of its location
\end{enumerate}
\par
\vspace{0.25\baselineskip}
\noindent\textbf{47.} According to the passage, lighthouses can assist navigators in
\begin{enumerate}[label=(\Alph*),leftmargin=*,itemsep=0.2em]
\item identifying underwater obstructions
\item preparing their boats for advisories.
\item finding their positions in the open sea
\item ignoring their charts of inland channels
\end{enumerate}
\par
\vspace{0.25\baselineskip}
\noindent\textbf{48.} In line 13, the phrase "zero in on" is closest in meaning to
\begin{enumerate}[label=(\Alph*),leftmargin=*,itemsep=0.2em]
\item cast about for
\item disguise
\item pinpoint
\item point out
\end{enumerate}
\par
\vspace{0.25\baselineskip}
\noindent\textbf{49.} In line 15, the word "navigable" is closest in meaning to
\begin{enumerate}[label=(\Alph*),leftmargin=*,itemsep=0.2em]
\item nauseating
\item passable
\item pernicious
\item calamitous
\end{enumerate}
\par
\vspace{0.25\baselineskip}
\noindent\textbf{50.} What can be inferred from the last paragraph?
\begin{enumerate}[label=(\Alph*),leftmargin=*,itemsep=0.2em]
\item Lighthouses do not contain a living space for maintenance personnel.
\item Lighthouse beams are projected intermittently during nighttime.
\item Technological expertise is expected of the lighthouse maintenance personnel.
\item Lighthouse technology is outdated and should have been replaced.
\end{enumerate}
\par
\vspace{0.25\baselineskip}
\end{document}
